
\section{Glossario dei termini}
\begin{tblr}{h{3cm}h{2cm}m{10cm}h{3cm}} % - Ajustar esto
	\hline
	\textbf{Termine} & \textbf{Sinonimi} & \textbf{Definizione} & \textbf{Esempi d'uso} \\
	\hline
	% - Dobbiamo aggiornare questa tabella 
	Libretto & N/A & Il testo o "copione" di un'opera lirica.\\
	Opera & N/A & Un'opera lirica, uno spettacolo in cui si abbinano musica e canto lirico ad una messa in scena teatrale. \\
	Regia & N/A & I costumi, i fondali, la direzione coreografica e degli effetti speciali e, in generale, della narrazione visiva dell'opera. Periodicamente le opere ricevono una nuova regia, che viene realizzata da un \textit{regista}. \\
	Spettacolo & N/A & Uno spettacolo teatrale, comunemente la regia di un'opera o un concerto. \\
	Evento & N/A & La data e l'ora in cui un determinato spettacolo viene messo in scena. Una regia viene solitamente messa in scena più volte, mentre un concerto ha, di norma, una singola data. \\
	Sezione & N/A & Settore di Posti ... & L'Utente crea una \textbf{Sezione} \\
	sezione & N/A & Pagina centrale dell'app dove si può accedere a ... & L'Utente si trova nella \textbf{sezione} ... \\
	Cliente & N/A & L'aquirente di biglietti. \\
	Biglietteria & N/A & Un dipendente dell'ufficio biglietteria. \\
	Amministratore & N/A & Un gestore del sistema. \\
	Utente & N/A & Uno qualunque tra \textit{cliente}, \textit{biglietteria} e \textit{amministratore}. \\
	INFORMAZIONI & N/A & Una sezione dell'app dove si può accedere alle Opere, le Regie ed i Generi disponibili & L'Utente si trova nella sezione \textbf{INFORMAZIONI} \\
	SPETTACOLI & N/A & Una sezione dell'app dove si può gestire la creazione, modifica ed elimina degli Spettacoli e gli Eventi disponibili & \emph{Come nel caso precedente} \\
	PRENOTAZIONI & N/A & Una sezione dell'app dove si può gestire lo stato di pagamento e l'eliminazione delle Prenotazioni registrate & \emph{Come nel caso precedente} \\
	TEATRO & N/A & Una sezione dell'app dove si può gestire la creazione, modifica ed elimina delle Sezioni ed i Posti & \emph{Come nel caso precedente} \\
	\hline
\end{tblr}
\newpage